\documentclass{beamer}

\usepackage{amsmath,amsfonts,amssymb,bm}
\usepackage{quantikz}
\usepackage{braket}
\usepackage{physics}
\usepackage{graphicx}

\usetheme{Madrid}

\title{Quantum Neural Networks vs Classical Neural Networks}
\subtitle{A Physics-Oriented Graduate Lecture}
\author{}
\date{}

\begin{document}

\frame{\titlepage}

%================================================
\section{Introduction}
%================================================

\begin{frame}{Motivation}
\begin{itemize}
\item Classical neural networks are powerful function approximators
\item Quantum systems naturally live in exponentially large Hilbert spaces
\item Question: Can quantum systems enhance learning?
\end{itemize}

\medskip

\begin{block}{Core Idea}
Quantum neural networks (QNNs) use parameterized quantum circuits to perform learning tasks.
\end{block}
\end{frame}

%================================================
\section{Classical Neural Networks}
%================================================

\begin{frame}{Structure of a Neural Network}
\[
h^{(l+1)} = \sigma(W^{(l)} h^{(l)} + b^{(l)})
\]

\begin{itemize}
\item $W^{(l)}$: weights
\item $\sigma$: nonlinear activation
\end{itemize}
\end{frame}

%------------------------------------------------

\begin{frame}{Learning Problem}
Given data $(x_i,y_i)$:
\[
\min_\theta \sum_i (f(x_i;\theta) - y_i)^2
\]

Training via gradient descent:
\[
\theta \leftarrow \theta - \eta \nabla_\theta C
\]
\end{frame}

%================================================
\section{Quantum Neural Networks}
%================================================

\begin{frame}{Definition of a QNN}
\[
\ket{\psi(x,\theta)} = U(\theta)U(x)\ket{0}
\]

\[
f(x,\theta) = \bra{\psi} O \ket{\psi}
\]

\begin{itemize}
\item $U(x)$: data encoding
\item $U(\theta)$: trainable unitary
\item $O$: observable
\end{itemize}
\end{frame}

%------------------------------------------------

\begin{frame}{Basic QNN Circuit}
\[
\begin{quantikz}
\lstick{\ket{0}} & \gate{U(x)} & \gate{U(\theta)} & \meter{} \\
\lstick{\ket{0}} & \gate{U(x)} & \gate{U(\theta)} & \meter{}
\end{quantikz}
\]
\end{frame}

%================================================
\section{Feature Encoding}
%================================================

\begin{frame}{Data Encoding}
\[
U(x) = \exp\left(i \sum_i x_i Z_i\right)
\]

\begin{itemize}
\item maps classical data into quantum states
\end{itemize}
\end{frame}

%------------------------------------------------

\begin{frame}{Example Encoding Circuit}
\[
\begin{quantikz}
\lstick{\ket{0}} & \gate{R_z(x_1)} & \qw \\
\lstick{\ket{0}} & \gate{R_z(x_2)} & \qw
\end{quantikz}
\]
\end{frame}

%================================================
\section{Variational Ansatz}
%================================================

\begin{frame}{Parameterized Circuit}
\[
U(\theta) = \prod_i R_y(\theta_i)\prod_{i<j}\mathrm{CNOT}_{ij}
\]
\end{frame}

%------------------------------------------------

\begin{frame}{Example Ansatz Circuit}
\[
\begin{quantikz}
\lstick{\ket{0}} & \gate{R_y(\theta_1)} & \ctrl{1} & \qw \\
\lstick{\ket{0}} & \gate{R_y(\theta_2)} & \targ{} & \qw
\end{quantikz}
\]
\end{frame}

%================================================
\section{Measurement and Output}
%================================================

\begin{frame}{Observable Output}
\[
f(x,\theta) = \langle Z_1 \rangle
\]

\begin{itemize}
\item measurement introduces nonlinearity
\end{itemize}
\end{frame}

%================================================
\section{Training}
%================================================

\begin{frame}{Cost Function}
\[
C(\theta) = \sum_i (f(x_i,\theta) - y_i)^2
\]
\end{frame}

%------------------------------------------------

\begin{frame}{Parameter Shift Rule}
\[
\frac{\partial f}{\partial \theta}
=
\frac{1}{2}\left[f(\theta+\frac{\pi}{2}) - f(\theta-\frac{\pi}{2})\right]
\]

\begin{itemize}
\item exact gradient evaluation
\end{itemize}
\end{frame}

%================================================
\section{Comparison: Classical vs Quantum}
%================================================

\begin{frame}{Key Differences}
\begin{tabular}{l|l|l}
 & Classical NN & QNN \\
\hline
State space & $\mathbb{R}^d$ & Hilbert space \\
Nonlinearity & activation & measurement \\
Parameters & weights & gate angles \\
\end{tabular}
\end{frame}

%------------------------------------------------

\begin{frame}{Conceptual Mapping}
\begin{itemize}
\item weights $\leftrightarrow$ Hamiltonian parameters
\item activations $\leftrightarrow$ measurement
\item forward pass $\leftrightarrow$ unitary evolution
\end{itemize}
\end{frame}

%================================================
\section{Physics Interpretation}
%================================================

\begin{frame}{QNN as Quantum Dynamics}
\[
|\psi\rangle = e^{-iH(\theta)}|\psi(x)\rangle
\]

\begin{itemize}
\item learning = tuning Hamiltonian
\end{itemize}
\end{frame}

%------------------------------------------------

\begin{frame}{Pauli Expansion}
\[
H(\theta) = \sum_\alpha \theta_\alpha P_\alpha
\]

\begin{itemize}
\item $P_\alpha$: Pauli strings
\end{itemize}
\end{frame}

%------------------------------------------------

\begin{frame}{BCH Expansion}
\[
U^\dagger O U = O + i[H,O] + \cdots
\]

\begin{itemize}
\item generates correlations
\end{itemize}
\end{frame}

%================================================
\section{Expressivity}
%================================================

\begin{frame}{Entanglement}
\[
S(\rho_A) = -\mathrm{Tr}(\rho_A \log \rho_A)
\]
\end{frame}

%------------------------------------------------

\begin{frame}{Expressivity}
\begin{itemize}
\item entanglement increases model capacity
\item too much $\rightarrow$ barren plateaus
\end{itemize}
\end{frame}

%------------------------------------------------

\begin{frame}{Barren Plateaus}
\[
\mathrm{Var}(\nabla C) \sim e^{-n}
\]
\end{frame}

%================================================
\section{Quantum Fisher Information}
%================================================

\begin{frame}{Quantum Fisher Information}
\[
F_{ij} = 4\mathrm{Re}(\langle \partial_i \psi | \partial_j \psi\rangle)
\]
\end{frame}

%------------------------------------------------

\begin{frame}{Natural Gradient}
\[
\dot{\theta} = F^{-1}\nabla C
\]
\end{frame}

%================================================
\section{Example}
%================================================

\begin{frame}{Two-Qubit QNN Example}
\[
\begin{quantikz}
\lstick{\ket{0}} & \gate{R_z(x_1)} & \gate{R_y(\theta_1)} & \ctrl{1} & \qw \\
\lstick{\ket{0}} & \gate{R_z(x_2)} & \gate{R_y(\theta_2)} & \targ{} & \qw
\end{quantikz}
\]
\end{frame}

%================================================
\section{Summary}
%================================================

\begin{frame}{Summary}
\begin{itemize}
\item QNN = parameterized quantum circuit
\item combines quantum mechanics and machine learning
\item learning = optimization in Hilbert space
\item strong connections to many-body physics
\end{itemize}
\end{frame}

\end{document}


\documentclass{beamer}

\usepackage{amsmath,amsfonts,amssymb,bm}
\usepackage{quantikz}
\usepackage{braket}
\usepackage{physics}
\usepackage{graphicx}

\usetheme{Madrid}

\title{Quantum Neural Networks}
\subtitle{Variational Circuits, Geometry, and Many-Body Physics}
\author{}
\date{}

\begin{document}

\frame{\titlepage}

%================================================
\section{Overview}
%================================================

\begin{frame}{Conceptual Structure}
A quantum neural network is defined by
\[
\ket{\psi(x,\theta)} = U(\theta)U(x)\ket{0}
\]

Output:
\[
f(x,\theta) = \bra{\psi} O \ket{\psi}
\]

\medskip

Interpretation:
\begin{itemize}
\item feature map $U(x)$ prepares state
\item ansatz $U(\theta)$ evolves system
\item observable extracts information
\end{itemize}
\end{frame}

%================================================
\section{QNN as Many-Body Wavefunction}
%================================================

\begin{frame}{Wavefunction Interpretation}
\[
\ket{\psi(x,\theta)} = e^{-iH(\theta)}\ket{\psi(x)}
\]

\begin{itemize}
\item QNN = parameterized many-body state
\item learning = tuning Hamiltonian parameters
\end{itemize}
\end{frame}

%------------------------------------------------

\begin{frame}{Pauli Expansion}
\[
H(\theta) = \sum_\alpha \theta_\alpha P_\alpha
\]

\begin{itemize}
\item $P_\alpha$: Pauli strings
\item structure resembles spin Hamiltonians
\end{itemize}
\end{frame}

%------------------------------------------------

\begin{frame}{Example Ansatz Circuit}
\[
\begin{quantikz}
\lstick{\ket{0}} & \gate{R_y(\theta_1)} & \ctrl{1} & \qw \\
\lstick{\ket{0}} & \gate{R_y(\theta_2)} & \targ{} & \qw
\end{quantikz}
\]
\end{frame}

%================================================
\section{BCH Expansion and Effective Operators}
%================================================

\begin{frame}{Baker-Campbell-Hausdorff Expansion}
\[
U^\dagger O U = O + i[H,O] + \frac{i^2}{2}[H,[H,O]] + \cdots
\]

\medskip

This shows:
\begin{itemize}
\item QNN generates nested commutators
\item increasing circuit depth $\rightarrow$ higher-order correlations
\end{itemize}
\end{frame}

%------------------------------------------------

\begin{frame}{Interpretation}
\begin{itemize}
\item equivalent to many-body perturbation expansion
\item similar structure as coupled cluster theory
\end{itemize}
\end{frame}

%================================================
\section{Loss Landscape}
%================================================

\begin{frame}{Cost Function}
\[
C(\theta) = \sum_i (f(x_i,\theta) - y_i)^2
\]
\end{frame}

%------------------------------------------------

\begin{frame}{Gradient Structure}
\[
\partial_\theta f = i\bra{\psi}[H_\theta,O]\ket{\psi}
\]

\begin{itemize}
\item gradients determined by commutators
\end{itemize}
\end{frame}

%------------------------------------------------

\begin{frame}{Landscape Geometry}
\begin{itemize}
\item highly non-convex
\item many local minima
\item flat regions possible
\end{itemize}
\end{frame}

%================================================
\section{Barren Plateaus}
%================================================

\begin{frame}{Definition}
Barren plateaus occur when
\[
\mathrm{Var}(\nabla C) \sim e^{-n}
\]

\begin{itemize}
\item gradients vanish exponentially
\end{itemize}
\end{frame}

%------------------------------------------------

\begin{frame}{Origin}
\begin{itemize}
\item random circuits approximate unitary 2-designs
\item expectation values concentrate
\end{itemize}
\end{frame}

%------------------------------------------------

\begin{frame}{Avoiding Plateaus}
\begin{itemize}
\item shallow circuits
\item problem-inspired ansätze
\item local cost functions
\end{itemize}
\end{frame}

%================================================
\section{Quantum Fisher Information}
%================================================

\begin{frame}{Definition}
\[
F_{ij} = 4\,\mathrm{Re}(\langle \partial_i \psi | \partial_j \psi\rangle)
\]
\end{frame}

%------------------------------------------------

\begin{frame}{Metric Interpretation}
\[
ds^2 = \sum_{ij} F_{ij} d\theta_i d\theta_j
\]

\begin{itemize}
\item defines geometry of parameter space
\end{itemize}
\end{frame}

%------------------------------------------------

\begin{frame}{Natural Gradient}
\[
\dot{\theta} = F^{-1}\nabla C
\]

\begin{itemize}
\item accounts for curvature of state space
\end{itemize}
\end{frame}

%================================================
\section{Time-Dependent Variational Principle}
%================================================

\begin{frame}{TDVP Principle}
\[
\delta \| (i\partial_t - H)|\psi\rangle \| = 0
\]
\end{frame}

%------------------------------------------------

\begin{frame}{Equations of Motion}
\[
\sum_j F_{ij}\dot{\theta}_j = C_i
\]

\[
C_i = \mathrm{Im}\langle \partial_i \psi | H | \psi\rangle
\]
\end{frame}

%------------------------------------------------

\begin{frame}{Interpretation}
\begin{itemize}
\item QNN training resembles classical dynamics
\item parameters evolve according to effective equations
\end{itemize}
\end{frame}

%================================================
\section{Quantum Control Perspective}
%================================================

\begin{frame}{QNN as Control Problem}
\[
U(\theta) = \mathcal{T} \exp\left(-i \int H(t) dt\right)
\]

\begin{itemize}
\item parameters $\theta$ correspond to control fields
\end{itemize}
\end{frame}

%------------------------------------------------

\begin{frame}{Learning as Control}
\begin{itemize}
\item optimize control parameters
\item achieve desired observable expectation
\end{itemize}
\end{frame}

%================================================
\section{Connection to Kernel Methods}
%================================================

\begin{frame}{Induced Kernel}
\[
K(x,x') = \langle \psi(x) | \psi(x') \rangle
\]
\end{frame}

%------------------------------------------------

\begin{frame}{Neural Tangent Kernel}
\[
K_{ij} = \frac{\partial f}{\partial \theta_i}
\frac{\partial f}{\partial \theta_j}
\]
\end{frame}

%================================================
\section{Expressivity}
%================================================

\begin{frame}{Entanglement}
\[
S(\rho_A) = -\mathrm{Tr}(\rho_A \log \rho_A)
\]
\end{frame}

%------------------------------------------------

\begin{frame}{Role of Entanglement}
\begin{itemize}
\item necessary for complex correlations
\item excessive entanglement leads to loss of structure
\end{itemize}
\end{frame}

%================================================
\section{Comparison with Classical Neural Networks}
%================================================

\begin{frame}{Comparison}
\begin{tabular}{l|l|l}
 & Classical NN & QNN \\
\hline
Dynamics & linear + nonlinear & unitary \\
Parameters & weights & Hamiltonian parameters \\
Geometry & Euclidean & Riemannian (QFI) \\
\end{tabular}
\end{frame}

%------------------------------------------------

\begin{frame}{Key Differences}
\begin{itemize}
\item classical networks rely on explicit nonlinearities
\item QNNs rely on measurement-induced nonlinearity
\item optimization geometry fundamentally different
\end{itemize}
\end{frame}

%================================================
\section{Circuits}
%================================================

\begin{frame}{Full QNN Circuit}
\[
\begin{quantikz}
\lstick{\ket{0}} & \gate{U(x)} & \gate{R_y(\theta_1)} & \ctrl{1} & \qw \\
\lstick{\ket{0}} & \gate{U(x)} & \gate{R_y(\theta_2)} & \targ{} & \qw
\end{quantikz}
\]
\end{frame}

%================================================
\section{Summary}
%================================================

\begin{frame}{Summary}
\begin{itemize}
\item QNN = variational quantum state
\item training = geometric optimization
\item strong connection to many-body physics
\item challenges: barren plateaus and noise
\end{itemize}
\end{frame}

\end{document}
